\documentclass[11pt]{article}
\usepackage[margin=1in]{geometry}
\usepackage[T1]{fontenc}
\usepackage[utf8]{inputenc}
\usepackage{lmodern}
\usepackage{microtype}
\usepackage{xcolor}
\usepackage{amsmath,amssymb,amsthm}
\usepackage{enumitem}
\usepackage{hyperref}
\usepackage{setspace}
\usepackage{titlesec}
\usepackage{parskip}

\hypersetup{colorlinks=true,linkcolor=black,urlcolor=blue,citecolor=black}
\setstretch{1.06}
\setlist[itemize]{topsep=4pt,itemsep=2pt,leftmargin=1.4em}
\titleformat{\section}{\large\bfseries}{\thesection.}{0.6em}{}
\titleformat{\subsection}{\normalsize\bfseries}{\thesubsection.}{0.5em}{}

\newtheorem{prop}{Proposition}
\newtheorem{claim}{Claim}

\newcommand{\ALCIRCC}{\ensuremath{\mathcal{ALCI}_{RCC5}}}
\newcommand{\Tri}{\mathrm{Tri}}
\newcommand{\TNbr}{\mathrm{TNbr}}
\newcommand{\Dzero}{D_0}
\newcommand{\comp}{\mathrm{comp}}

\begin{document}

\begin{center}
    {\LARGE\bfseries Technical Response to the Second Revision of\\[2pt]
    \emph{``A Tableau Calculus for ALCIRCC5 with Tri-Neighborhood Blocking''}}\\[10pt]
    {\large April 2026}\par
    \vspace{6pt}
    {\normalsize Review memorandum}
\end{center}

\vspace{8pt}

\begin{abstract}
This revision is a meaningful improvement over the earlier tableau draft. The new
Tri-neighborhood invariant is a better blocking candidate than plain triangle-type-set
equality, the manuscript now treats the same-map obstruction explicitly, and Section~8.5
is commendably candid about several weak points. Nevertheless, I do not think the main
decidability claim is proved yet.

Two issues remain decisive. First, the termination proof still lacks a valid global
argument bounding the total number of node creations. Second, the soundness argument
breaks already at the level of the initial domains in Definition~5.4: same-map pairs can
have empty domains even before T-filtering begins. This gives a concrete counterexample
to the initial non-emptiness claim in Lemma~5.5. Beyond that, the revised proof of
Lemma~5.5 remains partly heuristic and partly experimental.

My conclusion is therefore the same in substance as before, though for narrower reasons:
the paper contains real progress, but Theorems~4.1,~5.9, and~6.1 should not yet be
presented as established results.
\end{abstract}

\section{Overall assessment}

The manuscript addresses exactly the right open problem. Wessel's report leaves the
decidability of \ALCIRCC{} open, while already making clear that complete-graph
semantics destroy the usual tree-model intuition behind standard tableau blocking.
The present paper proposes \emph{Tri-neighborhood blocking} as a middle position
between two known extremes:
\begin{itemize}
    \item type-equality blocking, which terminates but is too weak for sound unraveling, and
    \item node-identity profile blocking, which supports sound replay but tends not to terminate.
\end{itemize}
That is a good and worthwhile idea.

This second revision also improves on the first tableau draft in several important ways.
The paper no longer relies on the simple map-based T-closed assignment that failed for
same-map pairs; instead it introduces disjunctive pair domains and tries to let T-filtering
plus arc-consistency resolve the obstruction. The strengthening from \(\Tri(x)=\Tri(y)\)
to the neighborhood condition \(\TNbr(x)=\TNbr(y)\) is also a genuine improvement.
In particular, I no longer regard my earlier objection about a purely first-person
triangle invariant as the main difficulty.

Even so, the present revision still falls short of a proof of decidability.

\section{The termination theorem is still not proved}

The revised proof of Theorem~4.1 is better than the previous one, because it explicitly
recognizes that a bound on the \emph{number of simultaneously active nodes} does not by
itself bound the \emph{total} number of created nodes. Unfortunately, the replacement
argument still does not close the gap.

The proof has two local ingredients:
\begin{itemize}
    \item each node can fire an existential demand only finitely many times, because once a
    specific witness exists, that demand stays satisfied; and
    \item each node is claimed to have only finitely many blocked/unblocked phases, because
    its signature changes only finitely often.
\end{itemize}
Both points may well be true. But the final inference is invalid: from bounded branching
and finitely many active periods \emph{per node}, the proof concludes by an appeal to
K\"onig's lemma that the creation tree is finite.

That does not follow. A rooted tree may be finitely branching and may give every node only
one active period, yet still be infinite: the infinite unary chain is the simplest counterexample
to that inference pattern. To deduce finiteness, one needs an additional \emph{global}
well-founded measure, for example a uniform bound on branch length or a birth-event measure
that decreases across every unblocked expansion. The manuscript does not supply such a
measure.

Section~8.5 itself essentially confirms this: it says that the interaction between label
changes and Tri-reconfiguration in the pre-stable regime has not been fully controlled.
That is not just a missing closed form for the complexity bound; it is the missing final
step in the correctness of Theorem~4.1.

\section{A concrete defect in the initial-domain construction}

The sharper problem in the soundness part is Definition~5.4. The initial domain for a
non-adjacent pair \((d_1,d_2)\) in the unraveling is defined as
\[
\Dzero(d_1,d_2)=\{R\in NR^- : (tp(d_1),R,tp(d_2))\in P(G)\},
\]
where \(P(G)\) records only relations between \emph{distinct tableau nodes}.

That construction can already fail before any T-filtering takes place.

\begin{prop}
Definition~5.4 does not guarantee non-empty initial domains. In particular,
Lemma~5.5's initial non-emptiness claim is false as stated.
\end{prop}

\begin{proof}
Consider the concept
\[
C_0 := \exists PP.U \;\sqcap\; \exists DR.W,
\qquad\text{where } U := \exists DR.W.
\]
A saturated clash-free completion graph may contain three active nodes
\(r,u,w\) with labels
\[
L(r)=\{C_0\},\qquad L(u)=\{U\},\qquad L(w)=\{W\},
\]
and edges
\[
E(r,u)=PP,\qquad E(r,w)=DR,\qquad E(u,w)=DR.
\]
This is composition-consistent, since
\(DR\in \comp(PP,DR)\).
Node \(w\) witnesses the root's \(\exists DR.W\) demand, and the same node \(w\)
also witnesses \(u\)'s \(\exists DR.W\) demand.

Now unravel from the root copy \(d_0\). One child \(d_1\) is the direct DR-witness
with \(\mathrm{map}(d_1)=w\). A second child \(e\) is the PP-copy of \(u\), and \(e\)
has its own DR-child \(d_2\) with \(\mathrm{map}(d_2)=w\). Thus
\[
d_1\neq d_2,\qquad \mathrm{map}(d_1)=\mathrm{map}(d_2)=w,
\]
and \(d_1,d_2\) are non-adjacent in the unraveling.

Assume moreover that \(w\) is the only tableau node of type \(L(w)\) in the completion
graph. Then \(P(G)\) contains no triple of the form
\((L(w),R,L(w))\), because Definition~5.4 only records edges between distinct tableau
nodes. Therefore
\[
\Dzero(d_1,d_2)=\emptyset.
\]
So the initial-domain construction can already fail before T-filtering or arc-consistency is
applied.
\end{proof}

This observation also shows that the paper's discussion of same-map pairs is still too narrow.
Those pairs do not arise only from ``multiple copies of a blocked subtree''; they also arise
whenever the same active witness is reused by different parents in the completion graph.

\section{Lemma~5.5 remains heuristic even beyond the preceding defect}

Suppose the defect above were repaired, for example by adding special same-type domains.
Even then, Lemma~5.5 would remain unproved.

The revised argument says that the disjunctive domains can ``break mirror symmetry'' because
arc-consistency may narrow a problematic pair-type domain from something like
\(\{DR,PO,PPI\}\) to a safe subdomain such as \(\{DR,PPI\}\), and it points to
computational checks on graphs of size 8, 10, and 12 as evidence. That is useful evidence,
but it is not a theorem.

What is still missing is a general proof that for every pair \((d_1,d_2)\) and every third
node \(d_3\), the T-filtered domains retain enough support to survive the manuscript's very
strong arc-consistency rule. Section~8.5 says exactly this: a fully formal proof of
intra-subtree T-closure would strengthen the argument. I agree with that assessment.

Since Lemma~5.5 is the entry point to Lemma~5.8 and then to Theorem~5.9, the gap is not
local to a side remark. It blocks the soundness theorem.

\section{What the paper should claim at this stage}

One point in the paper's favor is its honesty. Section~8.5 explicitly identifies the
intra-subtree T-closure problem and the total-node-creation problem as places requiring
further scrutiny. That is good mathematical practice.

But those caveats are not minor. They concern the two central theorems needed for the main
result:
\begin{itemize}
    \item Theorem~4.1 (termination), and
    \item Theorem~5.9 (model construction from an open branch), hence also Theorem~6.1.
\end{itemize}
For that reason, the abstract, Section~8.1's summary table, and the conclusion still overstate
what has been proved. The paper does not yet establish decidability of \ALCIRCC{}.

A more accurate headline would be: the manuscript proposes a stronger blocking invariant,
gives promising computational evidence on the PO-incoherent benchmark, and isolates the
remaining global proof obligations more clearly than earlier drafts. That is meaningful progress,
but it is not the final theorem.

\section{Conclusion}

My assessment of this last revision is therefore:
\begin{itemize}
    \item The move from plain triangle-type-set blocking to Tri-neighborhood blocking is a real
    improvement.
    \item The paper's treatment of same-map pairs is conceptually better than the original
    map-based replay argument.
    \item Nevertheless, the main result is still not proved. The termination theorem lacks a
    valid global finiteness argument, and the soundness proof fails already at the initial-domain
    construction for same-map pairs.
\end{itemize}

So I would still not accept the claim that the decidability of \ALCIRCC{} has now been
settled. The manuscript is better, but it remains a substantial partial result rather than a
finished proof.

\vspace{10pt}
\hrule
\vspace{8pt}

\noindent\textbf{References}

\begin{enumerate}[label={[\arabic*]},leftmargin=2.2em]
    \item Claude and Michael Wessel. \emph{A Tableau Calculus for ALCIRCC5 with
    Tri-Neighborhood Blocking}. Second revision, April 2026.
    \item Michael Wessel. \emph{Qualitative Spatial Reasoning with the ALCIRCC Family --
    First Results and Unanswered Questions}. Technical Report FBI-HH-M-324/03,
    University of Hamburg, 2002/2003.
\end{enumerate}

\end{document}
